\documentclass[a4paper]{article}
\usepackage[margin=25mm]{geometry}
\usepackage{fancyhdr}

% Define a custom page style for the first page
\fancypagestyle{firstpage}{%
  \fancyhf{}  % clear all
  \fancyhead[L]{Full Name:
 \makebox[0.4\textwidth][l]{\dotfill}}
    % \leaders\hbox{.}\hskip 0.5\textwidth\hfill}
%  \kern 0.5em\makebox[0.5\textwidth][l]{\dotfill}}
  \renewcommand{\headrulewidth}{0pt}  % optional: remove the header rule
}

\let\oldmaketitle\maketitle
\renewcommand{\maketitle}{\oldmaketitle\thispagestyle{firstpage}}

\usepackage{polyglossia}
\setmainlanguage{english}
\setmainfont{TeX Gyre Pagella}
\setsansfont[Ligatures=TeX]{TeX Gyre Heros}
\usepackage{xeCJK}
\setCJKmainfont{Noto Sans CJK JP}
\setCJKsansfont{Noto Sans CJK JP}
\setCJKmonofont{Noto Sans CJK JP}
%
\title{Semantics: Metaphor Tutorial}
\author{Francis Bond \url{<bond@ieee.org>}}
\date{}%2011-08-15}
\usepackage{fontenc}
\usepackage{polyglossia}
\setmainlanguage{english}
\setmainfont{TeX Gyre Pagella}
\setsansfont[Ligatures=TeX]{TeX Gyre Heros}
\usepackage{xeCJK}
\setCJKmainfont{Noto Sans CJK JP}
\setCJKsansfont{Noto Sans CJK JP}
\setCJKmonofont{Noto Sans CJK JP}
%\newcommand{\ans}[1]{\hfill{#1}}
%\newcommand{\ans}[1]{}
%%% dynamically add answers
\input{answers}

\usepackage{xspace}
\newcommand{\ent}{\ensuremath{\Rightarrow}\xspace}
\newcommand{\nent}{\ensuremath{\not\Rightarrow}\xspace}
\newcommand{\Y}[1]{\textbf{#1}}

\usepackage{multicol}
%\Restriction{}
%\rightfooter{}
%\leftheader{}
%\rightheader{}
\usepackage{mygb4e}
\usepackage[e,j]{mtg2e}
\newcommand{\lex}[1]{\textbf{\textit{#1}}}
\newcommand{\lx}[1]{\textbf{\textit{#1}}}
\newcommand{\ix}{\ex\it}
%\newcommand{\eng}{\textit}
\newcommand{\con}[1]{\textsc{#1}}
\usepackage{url}
\usepackage[normalem]{ulem}
\newcommand{\ul}[1]{\uline{#1}}
\newcommand{\txx}[1]{\textbf{#1}}

\begin{document}
\maketitle
\begin{enumerate}

\item Lakoff and Johnson (1980) proposed the following image schemas for love:
  \begin{itemize}
  \item LOVE IS A JOURNEY
  \item LOVE IS A FORCE (electromagnetic, gravitational)
  \item LOVE IS WAR
  \end{itemize}
Organise the following metaphors into the above three schemas.
\begin{exe}
  \ex \textit{They lost their momentum}
  \abox{LOVE is a FORCE}
 \ex \textit{There were sparks between us}
  \abox{LOVE is a FORCE}
  \ex \textit{Look how far we've come}
  \abox{LOVE is a JOURNEY}
\ex \textit{We're at a crossroad}
  \abox{LOVE is a JOURNEY}
  \ex \textit{He overpowered her}
  \abox{LOVE is WAR}
\ex \textit{I could feel the electricity between them}
  \abox{LOVE is a FORCE}
\ex \textit{We'll just have to go our separate ways}
  \abox{LOVE is a JOURNEY}
\ex \textit{We can't turn back now}
  \abox{LOVE is a JOURNEY}
\ex \textit{His whole life revolves around her}
  \abox{LOVE is a FORCE}
\ex \textit{She is besieged by suitors}
  \abox{LOVE is WAR}
\ex \textit{They are uncontrollably attracted to each other}
  \abox{LOVE is a FORCE}
\ex \textit{He is known for his conquests}
  \abox{LOVE is WAR}
\end{exe}
Can you think of other metaphors that do not fit into the above three schemas?

\item Match the metonym with its type
  \begin{center}
  \begin{tabular}{lcl}
  a. The ham sandwich wants the bill.&~~~~~~~~~~& i. PART FOR WHOLE \\[1.5ex]
b. She drank the whole bottle.&&ii. ATTRIBUTE FOR PERSON \\[1.5ex]
c. Oxford published a new dictionary.&&iii. OBJECT FOR USER\\[1.5ex]
d. All hands on deck.&&iv. PLACE FOR INSTITUTION \\[1.5ex]
e. The university raised tuition.&&v. CONTAINER FOR CONTENT  \\[1.5ex]
f. The suits approved the merger.&&vi. INSTITUTION FOR PEOPLE \\[1.5ex]
  \end{tabular}
\end{center}
 \abox{a-iii, b-v, c-iv, d-i, e-vi, f-ii}
\item For any two languages that you know, discuss similarities and
  differences in conventionalized metaphors of body parts
  (e.g. \eng{hand of a watch}, \jpn[well-known (lit: face is wide)]{kao-ga hiroi}).
 \abox{\begin{itemize}
  \item \jpn{atama ga katai} (頭がかたい) ``stubborn --- head is hard''
    \\ HEAD is MIND; THOUGHT is OBJECT
  \item \jpn{kuchi ga karui} (口が軽い) ``loose lipped --- mouth is
    light''
    \\ can't keep a secret
    \\ MOUTH is VOICE; THOUGHT is OBJECT
  \item \jpn{kuchi ga katai} (口がかたい) ``close mouthed --  mouth is
    hard''
    \\ can keep a secret
    \\ MOUTH is VOICE; THOUGHT is OBJECT
  \end{itemize}}
 \end{enumerate}
\newpage
 \section*{Some Czech Examples}


\subsection*{Metonymy in Czech}

\begin{exe}
  \ex \eng{Dvojka u okna chce zaplatit.}
      \trans `The table of two by the window wants to pay.'
      \hfill (OBJECT FOR USER)

  \ex \eng{Vypila celou láhev.}
      \trans `She drank the whole bottle.'
      \hfill (CONTAINER FOR CONTENT)

  \ex \eng{Oxford vydal nový slovník.}
      \trans `Oxford published a new dictionary.'
      \hfill (PLACE FOR INSTITUTION)

  \ex \eng{Všechny ruce na palubu!}
      \trans `All hands on deck!'
      \hfill (PART FOR WHOLE)

  \ex \eng{Univerzita zvýšila školné.}
      \trans `The university raised tuition.'
      \hfill (INSTITUTION FOR PEOPLE)

  \ex \eng{Obleky schválily fúzi.}
      \trans `The suits approved the merger.'
      \hfill (ATTRIBUTE FOR PERSON)
\end{exe}
\begin{exe}
  \ex \eng{Čtu Kunderu.}
      \trans `I'm reading Kundera.'
      \hfill (PRODUCER FOR PRODUCT)

  \ex \eng{Poslouchám Mozarta.}
      \trans `I'm listening to Mozart.'
      \hfill (PRODUCER FOR PRODUCT)

  \ex \eng{Praha jedná s Bruselem.}
      \trans `Prague is negotiating with Brussels.'
      \hfill (PLACE FOR GOVERNMENT / INSTITUTION)

  \ex \eng{Potřebujeme víc hlav.}
      \trans `We need more heads.'
      \hfill (PART FOR WHOLE)
\end{exe}


\subsection*{Metaphor in Czech}
\begin{exe}
  \ex \eng{Nosím kontaktní čočky.}
      \trans `I wear contact lenses.'
      \hfill (\textit{čočka} = `lentil' / `lens')
\end{exe}
\begin{exe}
  % \ex \eng{Teď je mi to jasné.}
  %     \trans `Now it's clear to me.'

  \ex \eng{Osvětlil nám ten problém.} \hfill (KNOWLEDGE IS LIGHT)
      \trans `He shed light on that problem.'

  % \ex \eng{Vrhl na tu otázku nové světlo.}
  %     \trans `He cast new light on that question.'
% \end{exe}

% \subsubsection*{GOOD AS LIGHT / BAD AS DARK}

% \begin{exe}
  \ex \eng{Čeká nás světlá budoucnost.}  \hfill (GOOD IS LIGHT)
      \trans `A bright future awaits us.'

  \ex \eng{Jsou v ponuré náladě.}
      \trans `They are in a gloomy mood.'
% \end{exe}

% \subsubsection*{UP / DOWN metaphors}

% \begin{exe}
  \ex \eng{Je na dně.}  \hfill (GOOD IS UP)
      \trans `He's at rock bottom.'

  % \ex \eng{Ceny šly nahoru.}
  %     \trans `Prices went up.'

  % \ex \eng{Nálada šla dolů.}
  %     \trans `The mood went down.'

  \ex \eng{Morálka stoupla.}
      \trans `Morale rose.'
% \end{exe}

% \subsubsection*{UNDERSTANDING AS DEPTH}

% \begin{exe}
  \ex \eng{Jde do hloubky.}  \hfill (UNDERSTANDING IS DEPTH)
      \trans `He goes into depth.'

  \ex \eng{Má hluboké znalosti češtiny.}
      \trans `She has deep knowledge of Czech.'

  \ex \eng{To byl povrchní výklad.}
      \trans `That was a superficial explanation.'
\end{exe}

\end{document}



%%% Local Variables: 
%%% coding: utf-8
%%% mode: latex
%%% TeX-PDF-mode: t
%%% TeX-engine: xetex
%%% End: 
